\documentclass[10pt,a4paper]{amsart}
\usepackage[margin=19mm]{geometry}
\usepackage{amsmath,amssymb,mathtools}
\usepackage[scr=rsfs,cal=euler]{mathalfa}
\usepackage{xfrac}
\usepackage[unicode,hidelinks,bookmarks=false]{hyperref}
\newcommand{\ie}{i.e.\ }
\newcommand{\cf}{cf.\ }
\newcommand{\resp}{respectively\ }
\newcommand{\Subsection}{\subsection}
\providecommand{\nobreakdash}{\nobreak\hbox{-}}
\setlength{\emergencystretch}{2em}
\multlinegap0pt
\usepackage{amssymb}
\usepackage{mathtools}
\usepackage{xcolor}
\newtheorem{lemma}{Lemma}
\title{Explicit formulas for the Casimir eigenvalues of $SL(n,\mathbb{Z})$-Maass forms}
\author{Vishal Muthuvel}
\date{Source: arXiv:2605.16803v1 (CC BY 4.0)}
\begin{document}
\maketitle

\noindent\emph{Source and licence.} Explicit formulas for the Casimir eigenvalues of $SL(n,\mathbb{Z})$-Maass forms. Version arXiv:2605.16803v1 (CC BY 4.0), distributed by arXiv under the Creative Commons Attribution 4.0 International licence, http://creativecommons.org/licenses/by/4.0/, as recorded in the arXiv licence metadata for this exact version and shown on the abstract page https://arxiv.org/abs/2605.16803v1. Full licence notice: \texttt{licenses/arxiv-2605.16803-CC-BY-4.0.txt}.\par\smallskip

\noindent\textit{Editorial scope.} Counted unit: the article's lemma lemma:power-rule (source0.tex lines 642--648). The article's complete original proof of it is delivered with it (source0.tex lines 650--693, one \texttt{proof} environment of 44 lines). No mathematics is written, repaired or re-derived: the statement and the proof are the article's own lines, in the article's own notation, and no retained line is substituted or re-flowed.  No further source-local context is needed: every cross-reference of every retained line resolves inside the two delivered spans.  Nothing was omitted from any retained span, so the package carries no omission record.  No retained line contains a citation, so the wrapper carries no bibliography and no bibliography entry.  No retained line contains a figure environment, an image-inclusion command or drawing code, so no image is delivered and the package declares no asset dependency.  Editorial formatting only, in addition: an AMS \texttt{amsart} preamble that restates the article's own declarations and macros used by the retained lines, namely the environments \texttt{lemma} on separate counters, together with the standard packages those lines need.  The retained environments print in this document as Lemma 1 in order of appearance, whereas the article prints the same items with the numbering of its own full document.  The complete native source of the article as distributed in the arXiv e-print for this exact version is bundled unchanged as \texttt{sources/200748/source0.tex}.
\begin{lemma}\label{lemma:power-rule}
Let $\beta \in \mathbb{C}$, and let $f:\mathbb{R}^m\to\mathbb{C}$ be a complex-valued function in $m$ real variables $t_1,\dots,t_m$ that is at least $m$ times differentiable in a neighborhood of the origin $\mathbf{0}$. Then the power function $f^\beta:\mathbb{R}^m\to\mathbb{C}$ is also a complex-valued function in $t_1,\dots,t_m$ that is at least $m$ times differentiable in some neighborhood of $\mathbf{0}$. Furthermore, for any non-empty subset $S \subset \{1, \dots, m\}$, the mixed partial derivative of $f^\beta$ at the origin with respect to $t_j$ ($j\in S$) is 
\begin{equation}\label{eq:power-rule}
\frac{\partial^{|S|}(f^\beta)}{\prod_{j \in S}\partial t_j}(\mathbf{0}) \ = \ \sum_{k=1}^{|S|} \ (\beta)_k \cdot f^{\beta-k}(\mathbf{0})\sum_{\substack{\sqcup_{i=1}^k S_i = S, \\ S_i \neq \emptyset}} \ \prod_{i=1}^k \ \frac{\partial^{|S_i|}f}{\prod_{j\in S_i}\partial t_j}(\mathbf{0}),
\end{equation}
where the sum is over unordered partitions of $S$, and $(\beta)_k:=\beta(\beta-1)\cdots(\beta-k+1)$ denotes the falling factorial. 
\end{lemma} 

\begin{proof}
Again, we prove this lemma via induction on $|S|$. The base case $|S|=1$ corresponds to the familiar power rule for single-variable derivatives. Letting $S=\{j\}\subset\{1,\dots,m\}$, we have
\begin{equation}
    \frac{\partial (f^\beta)}{\partial t_j}(\mathbf{0}) \ = \ \beta f^{\beta-1}(\mathbf{0}) \ \frac{\partial f}{\partial t_j}(\mathbf{0}).
\end{equation}

For the inductive case, we assume that the lemma holds when $1\leq|S|=m'<m$. Take any $S\subset\{1,\dots,m\}$ such that $|S|=m'+1$, and any $j'\in S$. We assumed that the lemma holds that for $S\backslash\{j'\}$, so 
\begin{equation}
\frac{\partial^{|S|} (f^\beta)}{\prod_{j \in S\backslash\{j'\}}\partial t_j}(\mathbf{0}) \ = \ \sum_{k=1}^{|S|-1} \ (\beta)_k \cdot f^{\beta-k}(\mathbf{0})\sum_{\sqcup_{i=1}^k S_i = S\backslash\{j'\}} \ \prod_{i=1}^k \ \frac{\partial^{|S_i|}f}{\prod_{j\in S_i}\partial t_j}(\mathbf{0}).
\end{equation}
Differentiating both sides with respect to $t_{j'}$ and applying the familiar product and power rules for single-variable derivatives, we get
\begin{align}
\nonumber
& \frac{\partial^{|S|} (f^\beta)}{\prod_{j \in S}\partial t_j}(\mathbf{0}) \\
\nonumber
= \ & 
\sum_{k=1}^{|S|-1} \ (\beta)_k \ \frac{\partial}{\partial t_{j'}}\left(f^{\beta-k}(\mathbf{0})\sum_{\sqcup_{i=1}^k S_i = S\backslash\{j'\}} \prod_{i=1}^k \ \frac{\partial^{|S_i|}f}{\prod_{j\in S_i}\partial t_j}\right)(\mathbf{0})\\
\label{eq:ind-power-rule1}
= \ & 
\sum_{k=1}^{|S|-1} \ (\beta)_{k+1} \cdot f^{\beta-k-1}(\mathbf{0}) \ \sum_{\substack{\sqcup_{i=1}^{k+1} S_i = S, \\ S_{k+1}=\{j'\}}} \ \prod_{i=1}^{k+1} \ \frac{\partial^{|S_i|}f}{\prod_{j\in S_i}\partial t_j}(\mathbf{0}) \\
\label{eq:ind-power-rule2}
+ \ & \sum_{k=1}^{|S|-1} \ (\beta)_k \cdot f^{\beta-k}(\mathbf{0})\sum_{\sqcup_{i=1}^k S_i = S\backslash\{j'\}} \frac{\partial}{\partial t_{j'}}\left(\prod_{i=1}^k \ \frac{\partial^{|S_i|}f}{\prod_{j\in S_i}\partial t_j}\right)(\mathbf{0})
\end{align}
Changing the index of summation in \eqref{eq:ind-power-rule1} and applying the product rule to \eqref{eq:ind-power-rule2}, we get
\begin{align}
\nonumber
& \frac{\partial^{|S|} (f^\beta)}{\prod_{j \in S}\partial t_j}(\mathbf{0}) \\
\nonumber
= \ & 
\sum_{k=2}^{|S|} \ (\beta)_k \cdot f^{\beta-k}(\mathbf{0}) \ \sum_{\substack{\sqcup_{i=1}^{k} S_i = S, \\ S_{k}=\{j'\}}} \ \prod_{i=1}^{k} \ \frac{\partial^{|S_i|}f}{\prod_{j\in S_i}\partial t_j}(\mathbf{0}) \\
\nonumber
+ \ & \sum_{k=1}^{|S|-1} \ (\beta)_k \cdot f^{\beta-k}(\mathbf{0})\sum_{\sqcup_{i=1}^k S_i = S\backslash\{j'\}} \sum_{i=1}^k \left(\frac{\partial^{|S_i|+1}f}{\prod\limits_{j\in S_i\cup\{j'\}}\partial t_j}\prod_{i'\neq i} \ \frac{\partial^{|S_{i'}|}f}{\prod_{j\in S_{i'}}\partial t_j}\right)(\mathbf{0}) \\
\label{eq:ind-power-rule3}
= \ & 
\sum_{k=2}^{|S|} \ (\beta)_k \cdot f^{\beta-k}(\mathbf{0}) \ \sum_{\substack{\sqcup_{i=1}^{k} S_i = S, \\ S_{k}=\{j'\}}} \ \prod_{i=1}^{k} \ \frac{\partial^{|S_i|}f}{\prod_{j\in S_i}\partial t_j}(\mathbf{0}) \\
\label{eq:ind-power-rule4}
+ \ & \sum_{k=1}^{|S|-1} \ (\beta)_k \cdot f^{\beta-k}(\mathbf{0})\sum_{\sqcup_{i=1}^k S_i = S} \ \prod_{i=1}^k \ \frac{\partial^{|S_{i}|}f}{\prod_{j\in S_{i}}\partial t_j}(\mathbf{0}).
\end{align}
The above equality can be justified using the same recursive argument for $k$-partitions of $S$ as in the proof of the previous lemma. Furthermore, note that the first sum \eqref{eq:ind-power-rule3} corresponds to (non-empty) partitions $\sqcup_{i=1}^k S_i$ of $S$ in which one of the partition elements is $S_k=\{j'\}$; and the second sum corresponds to partitions of $S$ in which none of the partition elements equal $\{j'\}$, i.e., the partition element containing $j'$ contains at least one other member of $S$ as well. Together, these sums account for all the partitions of $S$, so
\begin{align}
\frac{\partial^{|S|} (f^\beta)}{\prod_{j \in A}\partial t_j}(\mathbf{0}) \ = \ & \sum_{k=1}^{|S|} \ (\beta)_{k} \cdot f^{\beta-k}(\mathbf{0}) \ \sum_{\sqcup_{i=1}^{k+1} S_i = S} \ \prod_{i=1}^{k} \ \frac{\partial^{|S_i|}f}{\prod_{j\in S_i}\partial t_j}(\mathbf{0}).
\end{align}  
This completes the induction. 
\end{proof}

\end{document}
